\documentclass{beamer}
\usepackage{graphicx}
\usepackage{xcolor}
\usepackage{tikz}

% 自定义颜色
\definecolor{purpleblue}{RGB}{80, 80, 180}
\definecolor{lightgray}{RGB}{230, 230, 230}
\definecolor{novelPurple}{RGB}{98, 54, 255} 
\definecolor{darkgray}{RGB}{50, 50, 50}

% 主题设置
\usetheme{Madrid}  
\usecolortheme{dolphin} 

% 设置字体（移除 fontspec，确保 pdfLaTeX 兼容）
\setbeamerfont{title}{size=\Large,series=\bfseries}
\setbeamerfont{frametitle}{size=\LARGE,series=\bfseries}
\setbeamercolor{frametitle}{fg=white, bg=novelPurple}
\setbeamercolor{title}{fg=white, bg=purpleblue}

% 封面信息
\title[SAT Hard Question]{\textbf{SAT Hard Question 12}}
\author{\textbf{Jack Zeng}}
\institute{\textcolor{white}{\textbf{Novel Prep}}}
\date{\today} 

\begin{document}

% --- Title Slide ---
\begin{frame}
    \titlepage
    \vfill
    \begin{flushright}
        \textcolor{purpleblue}{\Huge \textbf{Novel Prep}}
    \end{flushright}
\end{frame}


\begin{frame}
    \begin{tikzpicture}[remember picture, overlay]
        % 设置浅色渐变背景
        \shade[top color=softPurple, bottom color=softGray] 
            (current page.north west) rectangle (current page.south east);
        
        % 添加高斯模糊光晕
        \fill[white, opacity=0.3] (current page.center) circle (4cm);
        
        % 主标题
        \node[align=center] at (current page.center) {
            {\Huge\bfseries\textcolor{purpleblue}{Novel Prep}} \\[0.5cm]
            {\Large\itshape\textcolor{lightText}{Excellence in SAT Prep}}
        };

        % 添加柔和光点（点缀）
        \foreach \x in {1,2,3,4,5} {
            \fill[purpleblue, opacity=0.2] (rand*10-5, rand*7-3.5) circle (0.1+\x*0.05);
        }

        ;
    \end{tikzpicture}
\end{frame}




\begin{frame}{\textbf{Question: Solving a Quadratic Equation}}

In the given equation, \( a \) and \( b \) are non-zero constants.  
Which of the following expressions gives a solution to the equation?

\[
7ax^2 + (7b - 12a)x - 12b = 0
\]

\bigskip

\textbf{Choices:}
\begin{itemize}
    \item[A.] \(- \dfrac{a}{b} \)
    \item[B.] \(- \dfrac{b}{a} \)
    \item[C.] \( \dfrac{a}{b} \)
    \item[D.] \( \dfrac{b}{a} \)
\end{itemize}

\end{frame}


\begin{frame}{\textbf{Answer Explanation: Solving a Quadratic Equation}}



\textbf{Correct Answer: B.} \( \boxed{-\dfrac{b}{a}} \)

\end{frame}


\begin{frame}{\textbf{Question: Equation with Infinitely Many Solutions}}

In the given equation, \( a \) and \( b \) are constants. The equation has infinitely many solutions. What is the value of \( b \)?

\[
\frac{ax + b}{3} = \frac{6}{17}x - \frac{13}{12}
\]

\textbf{Choices:}
\begin{itemize}
    \item[A.] \( -\frac{18}{17} \)
    \item[B.] \( \frac{18}{17} \)
    \item[C.] \( -\frac{13}{4} \)
    \item[D.] \( \frac{13}{4} \)
\end{itemize}

\end{frame}



\begin{frame}{\textbf{Answer Explanation: Equation with Infinitely Many Solutions}}

For two expressions to be equal for all values of \( x \), their corresponding coefficients must be equal.

\[
\frac{ax + b}{3} = \frac{6}{17}x - \frac{13}{12}
\]

Step 1: Match coefficients of \( x \):

\[
\frac{a}{3} = \frac{6}{17} \Rightarrow a = \frac{18}{17}
\]

Step 2: Match constants:

\[
\frac{b}{3} = -\frac{13}{12} \Rightarrow b = -\frac{13}{4}
\]

\textbf{Correct Answer: C.} \( \boxed{-\frac{13}{4}} \)

\end{frame}



\begin{frame}{\textbf{Question: Parabola and Intercepts}}

In the \( xy \)-plane, a parabola has \( x \)-intercepts at \( (-17, 0) \) and \( (37, 0) \).

If the equation of the parabola is written as
\[
y = ax^2 + bx + c,
\]
where \( a, b, c \) are constants, which of the following expressions gives \( c \) in terms of \( a \)?

\textbf{Choices:}
\begin{itemize}
    \item[A.] \( c = -1080a \)
    \item[B.] \( c = -629a \)
    \item[C.] \( c = \dfrac{a}{34} \)
    \item[D.] \( c = 51a \)
\end{itemize}

\end{frame}


\begin{frame}{\textbf{Answer Explanation: Parabola and Intercepts}}

\textbf{Solution:}

Given that the parabola has \( x \)-intercepts at \( -17 \) and \( 37 \), we can write the equation in factored form:
\[
y = a(x + 17)(x - 37)
\]

Expand the expression:
\[
y = a(x^2 - 20x - 629)
\]

Now match the standard form \( y = ax^2 + bx + c \). We identify:
\[
c = -629a
\]

\textbf{Correct Answer: B.} \( \boxed{c = -629a} \)

\end{frame}


\begin{frame}{\textbf{Question: Finding the Ratio of Triangle Segment Lengths}}

In triangle \( ABC \), the measure of angle \( B \) is \( 90^\circ \) and \( BD \) is an altitude of the triangle.  
The length of \( AB = 15 \), and the length of \( AC \) is 23 greater than the length of \( AB \).  

\medskip

What is the value of \( \dfrac{BC}{BD} \)?

\begin{itemize}
    \item[A.] \( \dfrac{15}{38} \)
    \item[B.] \( \dfrac{15}{23} \)
    \item[C.] \( \dfrac{23}{15} \)
    \item[D.] \( \dfrac{38}{15} \)
\end{itemize}

\end{frame}



\begin{frame}{\textbf{Answer Explanation: Using Right Triangle Properties}}
\small

\textbf{Step-by-step:}

\begin{itemize}
    \item \( AB = 15 \), \( AC = AB + 23 = 38 \)
    \item Use the Pythagorean Theorem to find \( BC \):
    \[
    AB^2 + BC^2 = AC^2 \Rightarrow 15^2 + BC^2 = 38^2
    \]
    \[
    225 + BC^2 = 1444 \Rightarrow BC^2 = 1219 \Rightarrow BC = \sqrt{1219}
    \]
    \item Area using base and height:
    \[
    \text{Area} = \frac{1}{2} \cdot AB \cdot BC = \frac{1}{2} \cdot AC \cdot BD
    \Rightarrow 15 \cdot BC = 38 \cdot BD
    \Rightarrow \frac{BC}{BD} = \frac{38}{15}
    \]

    \item \textbf{Correct Answer:} \( \boxed{\dfrac{38}{15}} \)
\end{itemize}

\end{frame}


\begin{frame}{\textbf{Question: Investment Account Model}}
\footnotesize

A certain investment account offers a special interest rate for the first 4 months the account is open followed by a lower interest rate for the remainder of the time the account is open.  

Akeyla opened one of these accounts with an original account balance of \$700.00 and did not make any other deposits or withdrawals.  
\begin{itemize}
    \item 4 months after Akeyla opened the account, the balance had increased by 0.6\% of the original balance.  
    \item 6 months after opening, the balance had increased by an additional 0.1\% of the balance at the end of the first 4 months.  
    \item Every 2 months after the first 6 months, the balance increased by an additional 0.1\% of the balance 2 months before.
\end{itemize}

Which of the following equations could represent the account balance \( B(x) \), in dollars, \( x \) months after the account was opened, where \( x \geq 4 \)?

\begin{itemize}
    \item[A.] \( B(x) = 704.20(1.001)^{\frac{x-4}{2}} \)
    \item[B.] \( B(x) = 700.00(1.001)^{\frac{x-4}{2}} \)
    \item[C.] \( B(x) = 704.20(1.001)^{2(x-4)} \)
    \item[D.] \( B(x) = 700.00(1.001)^{2(x-4)} \)
\end{itemize}

\end{frame}


\begin{frame}{\textbf{Answer Explanation: Investment Growth Equation}}

\textbf{Solution:}
\begin{itemize}
    \item Initial balance: \$700.00
    \item After 4 months, increased by 0.6\%: 
    \[
    700 + 0.006 \cdot 700 = 700 + 4.20 = 704.20
    \]
    \item From that point on, the account grows by 0.1\% every 2 months. So:
    \[
    \text{Multiplier every 2 months} = 1.001
    \]
    \item Since the growth starts from month 4, the number of 2-month periods since month 4 is:
    \[
    \frac{x - 4}{2}
    \]
    \item Thus, the model is:
    \[
    B(x) = 704.20(1.001)^{\frac{x-4}{2}}
    \]
\end{itemize}

\textbf{Correct Answer:} \( \boxed{\text{A}} \)

\end{frame}


\begin{frame}{\textbf{Question: Unit Conversion – Area Growth Rate}}

The area of a rectangular region is increasing at a rate of 180 square centimeters per hour.  

What is this rate, in square meters per minute?

\begin{itemize}
    \item[A.] 3
    \item[B.] 0.3
    \item[C.] 0.03
    \item[D.] 0.0003
\end{itemize}

\end{frame}
\begin{frame}{\textbf{Answer Explanation: Unit Conversion – Area Growth Rate}}

\textbf{Step-by-step solution:}

\begin{itemize}
    \item We are given a rate of \( 180 \ \text{cm}^2/\text{hour} \).
    \item First, convert square centimeters to square meters:
    \[
    1 \ \text{m}^2 = 10{,}000 \ \text{cm}^2 \quad \Rightarrow \quad \frac{180}{10{,}000} = 0.018 \ \text{m}^2/\text{hour}
    \]
    \item Then, convert hours to minutes:
    \[
    \frac{0.018}{60} = 0.0003 \ \text{m}^2/\text{min}
    \]
\end{itemize}

\textbf{Correct Answer:} \( \boxed{\text{D}} \)

\end{frame}

\begin{frame}{\textbf{Question: Exponential Function Evaluation}\textbf{Have soem Problems
}}

The exponential function \( f(x) \) is defined by 
\[
f(x) = ab^x
\]
where \( a \) and \( b \) are positive constants.  

If 
\textbf{Given:}
\[
f(n - 1) = f(n) + \frac{82}{100} f(n - 1)
\]  
where \( n \) is a constant, what is the value of \( b \)?

\begin{itemize}
    \item[A.] 1
    \item[B.] 0.82
    \item[C.] 0.18
    \item[D.] 10
\end{itemize}

\end{frame}
\begin{frame}

\textbf{Solve for } \( b \):
\[
b =\Rightarrow b = 0.18
\]

\textbf{Correct Answer:} \( \boxed{\text{C}} \)

\end{frame}

\begin{frame}{\textbf{Question: Exponential Function and Growth Rate}}

The function \( f \) is defined by the given equation, where \( k \) is a constant.  
The value of \( f(x) \) increases by 83\% for every increase of \( x \) by 1.

\[
f(x) = k(1.83)^x
\]

For which of the following functions, where \( k \) is a constant, does the value of \( g(z) \) increase by 83\% for every increase of \( x \) by \( \frac{1}{4} \)?

\begin{itemize}
    \item[A.] \( g(z) = k(1.83^{x})^{\frac{1}{4}} \)
    \item[B.] \( g(z) = k(1.83)^{x+\frac{1}{4}} \)
    \item[C.] \( g(z) = k(1.83^x)^4 \)
    \item[D.] \( g(z) = k(1.83)^{x - 4} \)
\end{itemize}

\end{frame}


\begin{frame}{\textbf{Question: Surface Area of Square Pyramids}}
A right square pyramid has a height of 11 centimeters (cm). A second right square pyramid has a height of 22 centimeters. The area of the bases of each of the two pyramids is 100 cm\(^2\).  
\vspace{0.3cm}

Which of the following is closest to the difference in the surface area of the second pyramid and the surface area of the first pyramid, in cm\(^2\)?

\begin{itemize}
    \item[(A)] 209
    \item[(B)] 220
    \item[(C)] 367
    \item[(D)] 551
\end{itemize}
\end{frame}

\begin{frame}{\textbf{Answer Explanation:  (Part 1)}}
\textbf{Step 1:} Surface area of a square pyramid is given by:  
\[
\text{Surface Area} = \text{Base Area} + \text{Lateral Area}
\]  
The lateral area for a square pyramid is:  
\[
\text{Lateral Area} = 2s \cdot l
\]  
where \(s\) is the side of the square base and \(l\) is the slant height.

\pause

\textbf{Step 2:} Base area is 100 cm\(^2\), so  
\[
s = \sqrt{100} = 10
\]

\textbf{Step 3:} Use Pythagoras’ theorem to find the slant height:
\[
l = \sqrt{\left(\frac{10}{2}\right)^2 + h^2} = \sqrt{25 + h^2}
\]

\textbf{First pyramid (\(h = 11\)):}  
\[
l_1 = \sqrt{25 + 121} = \sqrt{146}
\]

\textbf{Second pyramid (\(h = 22\)):}  
\[
l_2 = \sqrt{25 + 484} = \sqrt{509}
\]
\end{frame}

\begin{frame}{\textbf{Answer Explanation: Surface Area of Square Pyramids (Part 2)}}
\textbf{Step 4:} Calculate lateral areas:

\[
L_1 = 2 \cdot 10 \cdot \sqrt{146} \approx 20 \cdot 12.08 = 241.6
\]
\[
L_2 = 2 \cdot 10 \cdot \sqrt{509} \approx 20 \cdot 22.56 = 451.2
\]

\pause

\textbf{Step 5:} Add base areas:

\[
\text{Total Surface Area}_1 = 100 + 241.6 = 341.6
\]
\[
\text{Total Surface Area}_2 = 100 + 451.2 = 551.2
\]

\textbf{Step 6:} Find the difference:

\[
\Delta \text{Area} = 551.2 - 341.6 = 209.6 \approx \boxed{209}
\]

\textbf{Final Answer:} \textbf{(A) 209}
\end{frame}





\end{document}